\section{Conclusión}
Esta práctica fue una experiencia bastante reveladora sobre lo que realmente significa la seguridad en redes. Lo que parecía algo muy lejano y de película, como hackear un servidor o descubrir contraseñas, resultó ser algo completamente accesible con las herramientas correctas y un poco de paciencia.

Aprendimos que una contraseña débil o predecible es básicamente dejar la puerta de tu casa abierta con un letrero que dice ``pasen''. El simple hecho de que nuestra contraseña estuviera en un diccionario público fue suficiente para que Hydra la encontrara en cuestión de minutos. Eso nos hizo reflexionar mucho sobre nuestros propios hábitos digitales.

También entendimos por qué herramientas como nmap son tan poderosas y a la vez tan peligrosas. Con un solo comando pudimos ver qué puertos tenía abiertos un servidor en otro país, qué sistema operativo usaba y qué versión de SSH corría. Es información que en manos equivocadas puede ser el primer paso de un ataque real.

La parte de filtrar el diccionario con grep y awk nos enseñó algo muy valioso: en ciberseguridad, trabajar de forma inteligente es más importante que trabajar de forma bruta. Reducir millones de contraseñas a unas cuantas decenas de miles marcó la diferencia entre un ataque que tarda días y uno que termina en horas.

Finalmente, trabajar en equipo fue clave. Dividir el diccionario entre varios equipos y comunicarnos con el Comando Central para obtener pistas nos recordó que la seguridad, tanto del lado del atacante como del defensor, siempre es un esfuerzo colectivo.

En resumen, esta práctica nos bajó a tierra: la seguridad informática no es magia, es metodología, y protegerse tampoco lo es. Con políticas simples como usar llaves SSH, contraseñas largas y aleatorias, y herramientas como fail2ban, un servidor puede volverse muchísimo más difícil de comprometer. La diferencia entre ser vulnerable y estar protegido muchas veces es solo cuestión de buenos hábitos.
